\section{在固定网络上实现展开式}\label{sec:assembly}

第~\ref{sec:full} 节中的不同求和项最终可能具有不同的连线和查找表。为在一个网络上表示它们的和，我们把选取的求和项连同对其最终连线的赋值一起编码为一份运行记录。局部检验决定哪些运行记录有贡献，字典则提供相应的张量分量。规模上界取决于如何共享宿主图。所有块在宿主图同一位置使用同一地址，而一条连线上传输的比特可随各块独立变化。

\subsection{从运行记录到张量}

设 $K$ 为所用的 $k$ 的扩域，令 $h=\lceil\log_2 n\rceil$。固定 $P\le n^{O(1)}$ 个带编号的处理器，每个处理器具有至多 $h^c$ 比特的二进制辅助信息以及一个固定的外部端口有序列表。其局部记录包含猜测值以及至多 $h^c$ 条收发消息。每条消息由标识来源、目的地和端口的 $O(h)$ 比特键以及一个比特串组成。对于每个用到的键，若它在两端各出现恰好一次且所附比特串相同，则通信有效。

用一个固定的有限比特随机存取机程序，根据辅助信息和外部比特检查局部记录。它可以输出长度至多为 $U$ 的串，作为取值于 $K$ 的公共字典 $W$ 的参数。程序的决策只依赖于这些串。给定外部比特后，对每份被接受的运行记录，将 $W$ 返回的所有值相乘，再对所有被接受的记录求和，便得到运行记录张量的对应分量。记录格式是固定的，其中包括长度字段与填充方式。程序中所有随规模变化的数据均通过辅助信息输入。

每个处理器的计费量包括其确定性运行时间、辅助信息、外部比特、猜测值、带报头的消息和输出串。每份记录或每个条目至少计一个单位。该计费量由记录决定。一个局部地址或一个时间值占 $(\log h)^{O(1)}$ 比特。在如此长度的字上进行比特运算，至多给下文的界增加一个 $(\log h)^{O(1)}$ 因子。

\begin{lemma}[运行记录的实现]\label{lem:fixed-transcript}
设每份局部记录的计费量至多为 $V_{\max}=2^{(\log h)^{O(1)}}$。设 $V\le n^{O(1)}$ 既是每份被接受的运行记录的总计费量的上界，也是固定辅助信息与外部接口的总体积的上界，其中每个处理器至少计一个单位。若 $U\le h$，且定理~\ref{thm:sort} 的宽记录估计成立，则对每个固定的 $\tau>0$，运行记录张量都具有一个固定网络，其规模至多为
\[
                          h^\tau V+P h^{C_0}.
\]
各门都是三价门或元数为 $O(1+U+\log h)$ 的字典表。其标签占 $(\log h)^{O(1)}$ 比特。常数 $C_0$ 只依赖于 $\tau,c$、关于 $n$ 的多项式界以及宽记录估计。吸收误差因子前，$V_{\max}$ 中的指数只出现于 $(\log h)^{O(1)}$ 形式的因子内。
\end{lemma}

\begin{proof}
\emph{按记录长度分配空间。}
对从 $1$ 起直至不小于 $\min(V,V_{\max})$ 的最小二次幂为止的每个二次幂容量 $t$，预留 $N_t=\min(P,\lfloor 2V/t\rfloor)$ 个槽位。计费量为 $s$ 的处理器使用满足 $t\ge s$ 的最小容量。这一类中每份记录的计费量都大于 $t/2$，故 $N_t$ 个槽位已经足够。猜测相应的处理器标识符和记录。检验标识符的范围及计费量，要求被占用的槽位形成按标识符递增排列的一个前缀，并固定所有未用条目与填充位。所需的两个容量界为
\begin{equation}\label{eq:asm-seats}
 \sum_t tN_t=O\bigl(V(1+\log V_{\max})\bigr),\qquad
 \sum_t N_t=O\bigl(P(1+\log V_{\max})\bigr).
\end{equation}

\emph{检验计算及其字典参数。}
对于容量为 $t$ 的槽位，猜测至多 $t$ 步的确定性执行过程。依时间顺序检验指令、寄存器、计数器、地址及停机标志，停机后使用固定填充。为核验内存访问轨迹，把读、写和初始内容按地址和执行顺序排序。逐地址扫描，检验每次读取的都是最近的值，未使用单元从零开始。这就是 \citet[第~\ext{G}{app:trace}~节]{oaG} 的轨迹构造。对 $O(t)$ 次访问进行处理的比较网络，使用 $O(t\log^2(2t))$ 次对 $(\log h)^{O(1)}$ 比特字的比较。长输入、辅助信息、消息及输出串均存放于按地址访问的数组中，分段处理。局部费用为 $t(\log h)^{O(1)}$，每个槽位另加 $h^{O(1)}$，用于长标识符及记录分帧。

在每个槽位内，把输出串按其填充后可取的最小二次幂长度 $u$ 分组，空串计为长度一。在该类中预留 $\lfloor2t/u\rfloor$ 个调用，按发出顺序占用。检验程序将每个输出比特写入由其长度类别、序号和比特偏移量给出的地址。用同一内存轨迹检验，把这些数组与已装设的各次调用的输入逐一核对，包括启用标志、实际长度与零填充。每个地址占 $O(1+\log t+\log h)$ 比特。计入全部填充后的调用位置，会给检验费用增加 $t(\log h)^{O(1)}$。调用位置启用时，返回 $W$ 在指定串上的取值；未启用时，返回 $1$。相应的表由填充后的长度命名。其元数为 $O(1+U+\log h)$，每个槽位中所有表的总规模为 $O(t\log(2+U)\log(2+h))$。

\emph{匹配消息。}
固定 $w=\lceil h^{C_1}\rceil$，其中 $C_1>1$ 将根据宽记录排序的需要选定。把每条消息分成 $w$ 比特的分片，最后一片填充至满，空消息也预留一个分片。每份分片记录包含启用标志、键、方向、总长度、分片编号和有效载荷。给每个处理器增设一个固定发送方，其键属于预留的命名空间，负责提供该处理器的辅助信息、外部比特及各部分的长度。每个已占用的槽位必须接收对应于其所声明标识符的全部供给，即使供给为空也须如此。

将每条发出消息的来源和每条接收消息的目的地与槽位标识符核对，并检验分片方式、片段顺序、长度与填充。按启用标志、键、分片编号和方向对记录排序。在启用的部分，要求每个键和分片恰好对应一发一收，二者长度与有效载荷相同。一侧若有重复消息，其第零个分片就会冲突。特别地，每一份私有供给恰好匹配一个槽位，因此每个处理器恰好出现一次。对给定输入，稳定排序关系的输出唯一，完成上述检验后可对该输出求和消去。它在 $k$ 上的零一张量，扩域至 $K$ 后仍表示同一关系。

每个槽位预留 $O(t/w+h^c+1)$ 份分片记录。每个分片复制一次 $O(h)$ 比特的报头，因此构造这些记录的费用为 $(t+h^{O(1)})(\log h)^{O(1)}$。对固定发送方，按其实际供给和一个未满分片计费。由式~\eqref{eq:asm-seats}，包括消息边界处填充在内的预留比特总体积为
\[
                     V(\log h)^{O(1)}+P h^{O(1)} .
\]
将列表长度填充为二的幂。其长度是 $n$ 的多项式，故长度的对数为 $O(h)$。在定理~\ref{thm:sort} 中取容差 $\tau/3$，即可选定一个足够大的常数 $C_1$，使排序费用至多为上述体积的 $O(h^{\tau/3})$ 倍。随后的扫描具有线性规模。由式~\eqref{eq:asm-seats}，局部计算与字典表也满足同一体积界。将 $\log h$ 的幂吸收到 $h^\tau$ 中，并将各处理器的费用吸收到 $P h^{C_0}$ 中。

检验门和排序门均为 $k$ 上的三价门，每个字典门以其长度类别命名，名称的长度满足所述界。对于每份被接受的运行记录，槽位容量、标识符、顺序、轨迹、调用和分片均唯一确定。反过来，局部检验和通信检验能恢复一份被接受的记录。因此收缩已装设的关系，恰好得到 $K$ 上这些运行记录的加权和。
\end{proof}

\subsection{各轮的运行记录}

现在将引理~\ref{lem:fixed-transcript} 应用于第~\ref{sec:full} 节的展开。起初有 $S$ 个块，每块有 $T_0$ 个副本，故 $d=T_0S$。在 $j$ 轮中的每个阶段，一个分配了 $H$ 个位置的宿主图上有 $N_b\ge R$ 个块。给每个位置配备一个处理器，终末阶段也一样；每步另配 $O(H)$ 个处理器构成平衡汇总树。一个位置实例指某个块中的一个位置。

\paragraph{所有块共用一个地址。}
每个位置都存有其宿主图邻居的有序列表，以及一个按块顺序存放位置实例记录的数组。其辅助信息指定父位置和子位置的命名约定、整数参数、域的表示，以及外部端口的初始排序约定。初始位置的辅助信息还给出输入表的标签。外部比特通过固定端口按堆叠顺序进入。

在一步之内，宿主图所有顶点的度均有一个仅由固定的轮次设计决定的上界。因此，半径为 $O_D(1+\log r)$ 的球包含 $r^{O_D(1)}$ 个位置。在该半径内交换邻居列表，按标识符去除重复项，并将得到的名称排序。这样，每个可能的子位置和接触点都可由局部列表中的索引定位，每个位置需作 $h^{O(1)}$ 次公共操作。此后，发往接触点的一条消息就能用一个 $O(h)$ 比特的报头携带整个位置实例记录数组。数组内部只需局部索引。每个位置有 $h^{O(1)}$ 个阶段和声明。由式~\eqref{iter:labels}，每个位置处理这些数组的费用为 $N_b r^{O_D(1)}(\log h)^{O(1)}h^{O_D(a)}
=N_b h^{O_D(a)+o(1)}$。$O_D(a)$ 中的系数与 $a,b,\beta$ 无关。

对于宿主图共用的决策，每个位置计算各候选者评分中自己的局部贡献。沿平衡树对整数数组求和，并广播首个使评分最小的候选者。在高度 $i$ 处，初始为 $z$ 比特的条目之和至多有 $z+i$ 比特，而位置数以 $2^i$ 量级的因子缩减。由于 $\sum_{i\ge0}2^{-i}(z+i)^{O(1)}=O((1+z)^{O(1)})$，汇总保持上述比特费用界。在根处另允许 $2^{(\log h)^{O(1)}}$ 次操作以选取极小者。还可局部检验给定宿主图的贪心着色。按固定的标识符顺序，要求每个位置的颜色都取其较小编号的邻居尚未使用的最小颜色。按标识符归纳，被检验的赋色唯一。

\paragraph{检验宿主图的变更。}
一轮中的每项操作均可从上述邻居列表取得其接触点，并通过可加的评分取得公共参数。
\begin{enumerate}
\item 对聚落装包，尝试 $r$ 个槽位的乘子，将超载标记传播给祖先和邻居，并汇总各候选方案的聚落计数。每次传播都至少上升一个二进层级，故所需半径为 $O_D(1+\log r)$。相对区间端点、负载以及作二进缩放的计费量都具有简短编码。对局部评分求和也可选定用于完整条目配对的仿射着色。
\item 在局部集中步骤中，聚落区间及其原接触点给出每个球面的可能邻居。按已检验的颜色顺序，引理~\ref{lem:local} 中的移动选择与插值选择指定新表、取出的标量以及有序线束。该引理对描述长度的界足以支付选择与标签的费用。端口顺序确定露出的子列表。
\item 在节省连线的子轮中，来源聚落和容器决定已匹配单元所处的位置。每个实际使用的小构件记录一次求和项选择。对两种收缩分别求其保留下来的替换宽度之和，所有块均采用使该和最小的同一个选择。在合并重合端点之前得到的宽度已足以作此决策。随后由聚落的有界邻域确定新的接触点、端口顺序和表标签。
\item 重置时，计算每个位置对所有块的平均负载。平滑平均值的下限为 $FB/r$，传播的贡献每经过一条边至少衰减一半。初始负载至多为 $B$，故半径为 $O(1+\log r)$ 的邻域即可确定基区及其容量。

对公共混合器，令 $x_i$ 为块 $i\in\mathbb F_2^{\log_2 N_b}$ 内某位置的归一化负载，或某边的归一化宽度。当前缀长度为 $s$ 时，第~\ref{subsec:reset} 节碰撞估计量中依赖于候选矩阵 $A$ 的部分是下式的一个倍数：
\[
 \sum_{z\ne0}\boldsymbol{1}_{(Az)_{1:s}=0}
                     \sum_i x_i x_{i+z}.
\]
两个索引的输出前缀相同，当且仅当它们的差 $z$ 满足式中所列条件。对角项与中心化项关于 $A$ 为常数，所以选取极小者时可省去。以 $FB$ 乘以固定的二进负载分母，即可清除余下各 $x_i$ 的分母。经过固定缩放后，层级权重和溢出权重都是整数。因此，在任何特征下均可用整数运算作出该选择。

将沃尔什系数逐一平方，再作逆变换，即可同时算出各自相关量 $\sum_i x_i x_{i+z}$。这需要对 $(\log h)^{O(1)}$ 比特的局部整数作 $O(N_b\log N_b)$ 次运算。按索引之差和前缀长度汇总这些系数。在根处，花费至多 $2^{(\log h)^{O(1)}}$ 次比特运算即可枚举维数 $\log_2 N_b=O((\log h)^2)$ 下的所有可逆矩阵。第一个极小者满足平均论证给出的界。各位置据此置换自己的块数组。由于相交叶区间的长度比有界，且每个叶中的单元数有界，便可根据相邻记录确定终端分配、转子设置及新标签。
\item 对每个可能的站点与接触点，系线的最终分裂树规模都有一个依赖于 $D$ 的界。相邻位置决定哪些子列表处于启用状态。两个新子位置之间的一条边由其父位置之间一条长度有界的路径支撑。下一阶段中，子位置使用其自身宿主图的邻接关系，因此每次收集的邻域都是针对当前步骤的局部邻域。
\end{enumerate}

按阶段、端口和块的固定顺序声明各条消息。每次都依据前一阶段的数据和局部求和项选择来确定后续数组以及公共宿主图。局部恒等式中的每一项，包括权重为零的项，都只有一种规定的编码。因此，中间数组不会引入额外的重数。

\paragraph{独立比特与张量分量。}
按拆分与拼接所规定的顺序传送外部比特。对沿一个接触点发送的所有片段，来源与目的地共用一个处理器报头。各片段的数组位置与比特偏移量的长度为 $(\log h)^{O(1)}$。球面在其外部站点收集比特后，该站点在后续步骤中冻结。若各块重新排序，则该站点的列表也一并置换。收集邻域时传送的是描述。$ed$ 个外部比特中的每一个都只沿其指定路径移动，每次传送至多移动一次。因此 $O(\log h)$ 步的有效载荷处理工作量为 $ed(\log h)^{O(1)}$。

在终末阶段，于每条内部连线的指定端点处猜测一个比特，并按接触点和线束顺序将其发送至另一端。一项查表调用由站点描述及其有序端口比特组成。由式~\eqref{eq:round-terminal}，终末有效载荷的费用为 $L_j(\log h)^{O(1)}$。较早取出的标量在取出时发起调用，参数只有其简短描述。输入标量是固定的。若其乘积为 $c$，则用一个值为 $c^d$ 的零元表即可提供它们对整个展开式的贡献。命名 $K$ 中的这个元素需要 $O(\log h)$ 比特。给定各求和项的选择后，自由猜测的有效载荷恰好遍历最终连线的所有赋值。相应调用值的乘积就是该赋值的收缩权重。先对这些赋值求和，再对各求和项求和，得到在有序的 $ed$ 个外部比特上的输入张量幂。

\paragraph{固定网络的规模。}
在每个阶段，式~\eqref{eq:round-terminal} 给出 $HN_b\le mS h^{O_D(a)+o(1)}$。共有 $O(\log h)$ 步，且 $N_b\ge R$。因此处理器数与其总计费量分别满足
\begin{align}
 P&\le mS h^{O_D(a)+o(1)}/R,\nonumber\\
 V&\le P h^{C_2}+mS h^{O_D(a)+o(1)}
     +(L_j+ed)h^{o(1)}+2^{(\log h)^{O(1)}}.\label{eq:asm-volume}
\end{align}
前两项支付每个位置的公共工作以及每个位置实例的短记录的费用。第三项支付独立连线比特的费用，最后一项支付根处的计算费用。每份局部记录的计费量至多为 $2^{(\log h)^{O(1)}}$，因为 $N_b\le S$，$\log S=O((\log h)^2)$，且一个位置实例所含比特数仅为 $h$ 的多项式。指数 $C_2$ 以及用于控制声明数的指数，都可以独立于 $R=h^{M+o(1)}$ 中的 $M$ 固定下来。公共工作使用 $O(h)$ 比特的地址和 $r^{O_D(1)}$ 个候选接触点。对 $M$ 的依赖只出现在关于 $\log S$ 的多项式因子和作为加项的根计算费用中。

最长的字典调用占 $O_D(T_j)+h^{O_D(a)+o(1)}$ 比特。选 $a$ 充分小，使得对短记录工作量和标签界中的每个系数 $C$ 都有 $Ca+2a<\beta$，且对所有允许的 $b$ 都有 $b+(1+\delta)a<1$。由式~\eqref{eq:round-terminal} 得 $T_j=O(h^{b+(1+\delta/3)a})$，故各调用的长度均小于 $h$。在引理~\ref{lem:fixed-transcript} 中取 $\tau=a\delta/3$，再取 $M$ 充分大，以便同时吸收 $P h^{C_0}$ 和 $P h^{C_2+\tau}$。对于式~\eqref{eq:asm-volume} 中其余三项计费量，在乘以因子 $h^\tau$ 之前，每 $md$ 单位的费用分别为
\[
 \frac{mS h^{Ca+o(1)}}{md}=h^{-b+Ca+o(1)},\qquad
 \frac{ed}{md}\le h^{-2a},\qquad
 \frac{L_j}{md}=O(h^{-a}\log h).
\]
这里 $d=T_0S$ 且 $T_0=\Theta(h^b)$。作为加项的根计算费用为 $n^{o(1)}$，而 $m\ge n$。因此，当 $n$ 充分大时，固定网络的规模至多为 $md\,h^{-a+\tau+o(1)}\le md\,h^{-(1-\delta)a}$。包括启用标志字段和长度字段在内，其表的最大元数至多为 $h^{b+(1+\delta)a}$。字典表保留了其简短的长度类别标签，从而可再次应用引理~\ref{lem:macro}。至此即完成该引理的证明。
